\section{Experiments}
\label{sec:experiments}

\subsection{Setup}

Unless otherwise noted, experiments use \textbf{Qwen2.5-1.5B-Instruct}
\cite{qwen2.5} as the base LLM, extracting hidden states at layer 28 (the
final transformer block output). The nine stories per class come from the
pre-computed cache at
\texttt{line4/emotion\_vectors/qwen\_1.5B/stories\_*.json}. Bootstrap
confidence intervals use $1{,}000$ resamples. All datasets are loaded
via the \texttt{datasets} library at the versions current in
April~2026.

\subsection{Zero-Shot Sentiment Classification}

\paragraph{SST-2.}
We project the 872 SST-2 validation sentences onto the V-axis and threshold at
the median (yielding a zero-label binary classifier). We also report a
calibrated threshold tuned on a held-out 200-sample subset.

\begin{table}[h]
\centering
\caption{SST-2 zero-shot sentiment classification. A single unsupervised
direction extracted from 81 stories matches a supervised logistic
regression trained on 80--100 labeled SST-2 samples. 95\% bootstrap CIs
in brackets.}
\label{tab:sst2}
\begin{tabular}{lcc}
\toprule
Method & AUC & Accuracy \\
\midrule
Random direction (mean of 5) & 0.509 & 0.509 \\
SBERT MiniLM-L6 + same pipeline & 0.713 & 0.658 \\
Supervised LR, $N=10$ labels & 0.680 $\pm$ 0.058 & 0.608 $\pm$ 0.059 \\
Supervised LR, $N=100$ labels (L28 features) & 0.877 & 0.778 \\
Supervised LR, $N=5000$ labels (L28 features) & 0.923 & 0.848 \\
Supervised LR, $N=67{,}349$ (full SST-2 train, L28) & 0.951 & 0.878 \\
Direct LLM prompting (greedy) & -- & 0.928 \\
\rowcolor{blond}
\textbf{\ours (zero labels)} & \textbf{\sstauc{}} [0.844, 0.890] & \textbf{\sstacc{}} \\
\rowcolor{blond}
\textbf{\ours (200-sample threshold calib.)} & \sstauc{} & 79.17\% \\
\bottomrule
\end{tabular}
\end{table}

Our method at AUC $0.868$ sits between the supervised classifier with 50
and 100 labeled examples, while using zero labels and no gradient updates.

\paragraph{Four More Sentiment Benchmarks.}
We apply the same V-axis to IMDB (balanced 500+500), Yelp Polarity
(balanced 500+500), Rotten Tomatoes (1{,}066), and SemEval-2017 Twitter
sentiment (balanced 500+500).

\begin{table}[h]
\centering
\caption{The \emph{same} V-axis applied zero-shot to four additional
sentiment benchmarks (movies, restaurants, tweets, reviews). No
re-extraction, no labels.}
\label{tab:benchmarks}
\begin{tabular}{lccc}
\toprule
Benchmark & $n$ & AUC & Median-thresh Acc \\
\midrule
IMDB (balanced) & 1{,}000 & \imdbauc{} & \imdbacc{} \\
\textbf{Yelp Polarity (balanced)} & 1{,}000 & \textbf{\yelpauc{}} & \textbf{\yelpacc{}} \\
Rotten Tomatoes & 1{,}066 & \mrauc{} & 76.55\% \\
Twitter SemEval-17 & 1{,}000 & \twauc{} & 83.80\% \\
\bottomrule
\end{tabular}
\end{table}

All five English benchmarks use the \emph{identical} 81-story-centroid
V-axis. The same direction generalises across domains (Yelp restaurants,
IMDB / SST movies, Twitter), lengths (SST-2 sentences vs.\ IMDB
paragraphs), and registers (informal tweets vs.\ formal reviews).

\subsection{Word-Level Sentiment Lexicons}

\paragraph{Hu \& Liu 6{,}789 words.}
We extract the residual-stream embedding of each word in a minimal template
\texttt{"The word \textbackslash"\{w\}\textbackslash" means"}, then project onto
the V-axis.

\begin{table}[h]
\centering
\caption{Word-valence lexicon recovery. Pearson $|r|$ between our
V-axis projection and human-annotated valence scores.}
\label{tab:lexicon}
\begin{tabular}{lccc}
\toprule
Lexicon & $n$ & Pearson $|r|$ & Random baseline \\
\midrule
Hu \& Liu (binary +1/$-$1) & 6{,}789 & \textbf{\huliur{}} [0.74, 0.78] & 0.12 \\
Warriner \textit{et al.\@} (continuous 1--9) & 3{,}000 & \textbf{\warrinerr{}} & 0.07 \\
NRC 32-word fallback (curated) & 32 & 0.95 & 0.09 \\
\bottomrule
\end{tabular}
\end{table}

On both lexicons the signal-to-noise ratio vs.\ random projection is
$\geq6\times$. On Hu \& Liu the 95\% bootstrap confidence interval excludes
zero tightly ($[0.74, 0.78]$).

\paragraph{Continuous Sentence-Level Valence (EmoBank).}
EmoBank \cite{buechel2017emobank} provides 10{,}062 English sentences with
continuous VAD scores. Our V-axis projection correlates at
$|r|=\emobankr{}$ [0.485, 0.512] with ground-truth valence scores.
A supervised Ridge regression trained on the full feature space reaches
$|r|=0.66$; our method recovers $72\%$ of that signal with no labels.

\subsection{Fine-Grained Sentiment (SST-5)}

SST-5 provides 5-class ordinal valence (very negative \ldots very positive).
Our V-axis projection correlates at Pearson $r = 0.620$ (Spearman $0.627$)
with the 0--4 label on the 2{,}210-sample test set, confirming that the
axis captures continuous valence not just sign.

\subsection{Data Efficiency: How Few Stories Are Enough?}

\begin{figure}[t]
\centering
\includegraphics[width=0.7\textwidth]{nshot_curve}
\caption{Data-efficiency curve. \emph{Top}: SST-2 AUC as a function of the
number of stories per class (our method, zero labels). \emph{Bottom}:
supervised LR trained on $N$ labeled SST-2 samples. Our zero-label method
matches the supervised curve at $N=80$--$100$.}
\label{fig:nshot}
\end{figure}

A single story per class (9 total samples) already gives AUC 0.74 on SST-2.
Fifteen stories per class (135 total) match the supervised 5{,}000-label LR.

\begin{table}[h]
\centering
\caption{N-shot ablation on SST-2. Mean over 3 subsampling seeds.}
\label{tab:nshot}
\begin{tabular}{cc|cc}
\toprule
\multicolumn{2}{c}{\ours (zero labels)} & \multicolumn{2}{c}{Supervised LR ($N$ labels)} \\
$n$ per class & SST-2 AUC & $N$ labels & SST-2 AUC \\
\midrule
1 & 0.740 $\pm$ 0.043 & 10 & 0.680 $\pm$ 0.058 \\
3 & 0.738 $\pm$ 0.014 & 50 & 0.834 $\pm$ 0.035 \\
9 & 0.781 $\pm$ 0.029 & 100 & 0.878 $\pm$ 0.012 \\
15 & \textbf{0.831 $\pm$ 0.010} & 500 & 0.906 $\pm$ 0.003 \\
25 & 0.833 $\pm$ 0.008 & 2000 & 0.919 $\pm$ 0.001 \\
50 & 0.837 $\pm$ 0.000 & 5000 & 0.914 $\pm$ 0.000 \\
\bottomrule
\end{tabular}
\end{table}

At $n=1$ story per class (9 total LLM-generated samples with no labels),
we already beat the supervised classifier at $N=10$ labeled samples.

\subsection{The V-Shape: Valence Re-Crystallizes at the Top of the Network}
\label{sec:vshape}

\begin{figure}[t]
\centering
\includegraphics[width=0.9\textwidth]{vshape_qwen15B}
\caption{Per-layer SST-2 AUC of the V-axis extracted at each layer,
for Qwen2.5-1.5B (28 layers). Valence is present at layer~1 (embedding-level
signal), absent through layers 2--26 (the dilution zone), and
re-crystallizes at layers 27--28.}
\label{fig:vshape}
\end{figure}

\begin{table}[h]
\centering
\caption{Layer-wise SST-2 AUC for Qwen2.5-1.5B. The valence direction
re-crystallizes abruptly at layers 27--28.}
\label{tab:layerwise}
\begin{tabular}{lccccc}
\toprule
Layer & 1 & 14 & 26 & 27 & 28 \\
\midrule
SST-2 AUC & \textbf{0.752} & 0.525 & 0.532 & \textbf{0.836} & \textbf{0.837} \\
EmoBank V $|r|$ & 0.370 & 0.005 & 0.002 & 0.517 & 0.496 \\
\bottomrule
\end{tabular}
\end{table}

Cross-layer V-axis cosines are the mechanistic smoking gun: L1--L28
$= 0.29$, L14--L28 $= \textbf{0.07}$ (nearly orthogonal), L27--L28
$= 0.95$. The final-layer valence direction is \emph{not} carried
forward from earlier layers --- it is built de novo by the last one
or two transformer blocks, where the LM head must read it off.

\subsection{Cross-Family Universality}

\begin{table}[h]
\centering
\caption{Same method, five LLM families. EmoBank V $|r|$ at the final layer.}
\label{tab:crossmodel}
\begin{tabular}{lccc}
\toprule
Model & Family & $|V_L|$ & Recovery of Full+Ridge baseline \\
\midrule
Qwen2.5-0.5B-Instruct & Qwen & 0.485 & -- \\
\textbf{Qwen2.5-1.5B-Instruct} & Qwen & \textbf{0.475} & 72\% \\
Qwen2.5-7B-Instruct & Qwen & 0.455 & \textbf{77\%} \\
\midrule
Pythia-1.4b & GPT-NeoX & 0.461 & -- \\
Bloom-560m & BLOOM & 0.342 & -- \\
\bottomrule
\end{tabular}
\end{table}

All five models produce a usable late-layer valence direction. The
magnitude varies (Qwen strongest, BLOOM weakest) but every model shows the
same V-shape pattern.

\paragraph{Functional alignment across models.}
Different models have different hidden-state dimensionalities, so we
cannot directly compare their V-axis \emph{vectors}. Instead, we
project the same 81 stories onto each model's V-axis, producing a
per-story valence score for every model. After a simple sign alignment
(positive emotions on the positive side), we compute pairwise cosine
similarity and Spearman rank correlation of these 81-dimensional
per-story valence vectors.

\begin{table}[h]
\centering
\caption{Functional cross-model agreement on per-story valence
projections. Despite different architectures and sizes, all four models
rank the 81 stories nearly identically on valence.}
\label{tab:crossmodel_func}
\begin{tabular}{lcc}
\toprule
Model pair & Cosine sim & Spearman $\rho$ \\
\midrule
Qwen2.5-0.5B $\times$ Qwen2.5-1.5B & \textbf{0.980} & \textbf{0.977} \\
Qwen2.5-0.5B $\times$ Pythia-1.4b & 0.974 & 0.972 \\
Qwen2.5-0.5B $\times$ Bloom-560m & 0.944 & 0.935 \\
Qwen2.5-1.5B $\times$ Pythia-1.4b & \textbf{0.984} & \textbf{0.982} \\
Qwen2.5-1.5B $\times$ Bloom-560m & 0.944 & 0.932 \\
Pythia-1.4b $\times$ Bloom-560m & 0.937 & 0.924 \\
\midrule
\textbf{Mean (all 6 pairs)} & \textbf{0.961} & \textbf{0.954} \\
\bottomrule
\end{tabular}
\end{table}

Four models with different architectures, training corpora, and hidden
dimensionalities --- two Qwen scales, a GPT-NeoX, a BLOOM --- agree on
the \emph{ranking} of the 81 stories at mean pairwise cosine $0.961$ and
Spearman $\rho = 0.954$ (minimum $0.937$ for Pythia-1.4b vs.\
Bloom-560m). This is a stronger universality claim than the per-model
$|V_L|$ numbers in \cref{tab:crossmodel}: it says the \emph{same
per-story valence score} is recovered by every family we tested, up to
a sign. The late-layer valence direction looks like a convergent
solution of the next-token-prediction objective rather than an artifact
of a particular training pipeline.

\begin{figure}[t]
\centering
\includegraphics[width=0.9\textwidth]{vshape_crossmodel}
\caption{V-shape pattern on three model families. Top: Qwen2.5-1.5B (28
layers), middle: Pythia-1.4b (24 layers), bottom: Bloom-560m (24 layers).
All three show an early-layer signal (L1--L4), a middle-layer dilution
zone, and a final-layer re-crystallization.}
\label{fig:crossmodel}
\end{figure}

\subsection{Cross-Language Generalization}

The strongest claim of the paper is that an axis extracted from
\emph{English} stories alone classifies \emph{Chinese} text. We evaluate
two transfer protocols. \emph{Protocol 1 (zero-shot transfer).} Extract
the V-axis from the 81 English stories at Qwen2.5-1.5B L28. Score
held-out 64-sentence story splits in English, Spanish, French, German,
and Chinese --- no re-extraction, no re-fitting, no translation.
\emph{Protocol 2 (within-language cosines).} Extract a separate V-axis
per language (15 stories per emotion) and measure pairwise cosines.
Protocol 1 is the deployment-relevant question; Protocol 2 tests
whether the underlying direction is language-invariant at the
representation level.

\begin{figure}[t]
\centering
\includegraphics[width=0.85\textwidth]{fig_crosslang}
\caption{English V-axis transfers zero-shot to four non-English languages.
Bars show binary-sentiment AUC and median-threshold accuracy on held-out
64-sentence story splits; the English axis is reused with no translation
and no re-extraction. Chance line at $0.5$.}
\label{fig:crosslang}
\end{figure}

\begin{table}[h]
\centering
\caption{Cross-lingual transfer of the \emph{English} V-axis
(Qwen2.5-1.5B L28, zero labels). Chinese reaches perfect separation
(AUC $1.000$) without the model ever seeing Chinese during extraction.}
\label{tab:crosslang_transfer}
\begin{tabular}{lccc}
\toprule
Target language & $n$ & Transfer AUC & Transfer Acc. \\
\midrule
English (in-domain) & 64 & 0.990 & 93.75\% \\
Spanish & 64 & 0.949 & 90.63\% \\
French & 64 & 0.991 & 96.88\% \\
German & 64 & 0.972 & 90.63\% \\
\rowcolor{blond}
\textbf{Chinese} & 64 & \textbf{1.000} & \textbf{100.00\%} \\
\bottomrule
\end{tabular}
\end{table}

The English axis reaches AUC $\geq 0.949$ on every non-English test
split, including Chinese at $1.000$. Cross-lingual transfer requires
neither a multilingual encoder nor translation: the LLM's late-layer
residual stream places valence along a direction that is shared across
the languages it was trained on. For multilingual deployment this is a
single-axis probe that already works.

\begin{table}[h]
\centering
\caption{Within-language V-axis cosine similarity (Protocol 2).
Valence axes share roughly half their direction across languages; arousal
axes are essentially orthogonal.}
\label{tab:crosslang}
\begin{tabular}{lcc}
\toprule
Pair & V-axis cosine & A-axis cosine \\
\midrule
en vs.\ es & \textbf{0.466} & 0.186 \\
es vs.\ fr & \textbf{0.462} & $-$0.249 \\
de vs.\ en & 0.310 & 0.239 \\
de vs.\ es & 0.273 & 0.114 \\
en vs.\ fr & 0.162 & $-$0.097 \\
de vs.\ fr & 0.229 & 0.023 \\
\bottomrule
\end{tabular}
\end{table}

Within-language cosines are positive but modest ($[0.16, 0.47]$), indicating
that each language's 9-centroid geometry finds a direction partially but
not identically aligned with English's. The stronger transfer result
(Table~\ref{tab:crosslang_transfer}) is what matters for deployment: the
English axis is already a good cross-lingual probe in practice, even if
the underlying per-language directions differ by roughly $60^\circ$.

\subsection{Principal-Angle Analysis of the 9-Centroid Geometry}
\label{sec:principal_angle}

We verify that PC1 of the nine centroids is not a trivial
high-variance-axis artifact. The singular values of the mean-centered
centroid matrix at L28 are $(29.13, 15.21, 12.11, 10.17, 9.02, 6.61,
6.27, 4.96, \approx 0)$; PC1 captures $55.8\%$ of the centroid variance
and the ratio of the first two singular values is $1.91$, well separated.
The principal angle between PC1 and PC2 is $90^\circ$ (orthogonal by
construction); the principal angle between PC1 and a uniform random
direction in $\mathbb{R}^{1536}$ is $88.81^\circ \pm 0.83^\circ$ over
$100$ samples. The centroid PC1 is therefore a geometrically principled
direction: it is the unique single-direction summary that maximises
between-category spread, and it is not close to any random direction in
the high-dimensional hidden space.

\subsection{Model-Scale Stability}
\label{sec:scale_stability}

We extend the cross-family evaluation (Table~\ref{tab:crossmodel}) with
a scale sweep on the Qwen2.5 family (0.5B, 1.5B, 7B) at matched
extraction settings. All three scales produce a V-axis with EmoBank
V~$|r| \in [0.485, 0.494]$. The 1.5B and 7B models give essentially
identical valence correlation ($|r| = 0.494$ on both), indicating that
the recipe is not sensitive to model scale in the range
most commonly deployed.

\subsection{Valence/Arousal Asymmetry}
\label{sec:va_asymmetry}

\begin{table}[h]
\centering
\caption{Valence vs.\ arousal extraction success. Same method (PCA of 9 centroids)
applied to both. $|r|$ against EmoBank's continuous V and A scores.}
\label{tab:va_asymmetry}
\begin{tabular}{lcc}
\toprule
Layer & Valence $|r|$ & Arousal $|r|$ \\
\midrule
L1 & 0.370 & 0.012 \\
L14 & 0.005 & 0.013 \\
L27 & \textbf{0.517} & 0.007 \\
L28 & \textbf{0.496} & 0.044 \\
\bottomrule
\end{tabular}
\end{table}

Across every layer, every model family, and every continuous-regression
test, the valence direction is extractable and the arousal direction is not.
The maximum arousal $|r|$ we observed at any layer in any model is 0.16
(Warriner lexicon at L28), still well below the valence number.

\begin{figure}[t]
\centering
\includegraphics[width=0.85\textwidth]{va_asymmetry}
\caption{Valence vs.\ arousal layer-wise on Qwen2.5-1.5B. Valence (green)
has early-layer signal at L1, a middle dilution zone, and a sharp
re-crystallization at L27--L28 reaching $|r|=0.50$. Arousal (red) never
exceeds $|r|=0.05$ at any layer --- the asymmetry is structural, not a
small-effect question.}
\label{fig:va_asym}
\end{figure}

\subsection{PCA-of-Centroids Beats the Standard Contrastive Axis}
\label{sec:contrastive}

The dominant representation-engineering recipe
\cite{zou2023representation, rimsky2024steering, arditi2024refusal}
extracts a concept direction as the mean-difference of activations on
operator-specified contrastive pairs. We match it on the same features.
Re-partitioning our stories into a positive pool (\texttt{amusement,
inspiration, joy, tenderness}, 80 stories) and a negative pool
(\texttt{anger, fear, disgust, sadness}, 80 stories), we compute the
mean-difference axis at L28 and evaluate on the same 872 SST-2
validation sentences.

\begin{table}[h]
\centering
\caption{PCA-of-nine-centroids beats the standard mean-difference
contrastive axis by $+0.115$ AUC on matched features. The operator
specifies neither a positive/negative split nor a pairing --- PCA
discovers the polarity direction automatically.}
\label{tab:contrastive}
\begin{tabular}{lc}
\toprule
Axis extraction method & SST-2 val AUC \\
\midrule
Mean-diff contrastive (pos\,$-$\,neg, 80+80 stories) & 0.753 \\
\rowcolor{blond}
\textbf{PCA of 9 centroids (ours)} & \textbf{0.868} \\
\midrule
\emph{Absolute improvement} & \textbf{$+$0.115 AUC} \\
\bottomrule
\end{tabular}
\end{table}

The two axes have absolute cosine $0.527$ --- substantially aligned but
not identical. PCA of nine class-centroids improves SST-2 AUC by
$+0.115$ over the contrastive mean-difference baseline on matched
features, \emph{and} removes the requirement that the operator specify
a binary split. The same PCA axis transfers cross-lingually
(Table~\ref{tab:crosslang_transfer}) without needing a per-language
contrastive-pair specification. Taken together with the counterfactual
ablation, this is a direct methodological result: the $K$-centroid PCA
geometry extracts a stronger polarity direction than the binarised
contrastive alternative on the same hidden states, with strictly less
operator input.

\subsection{Counterfactual: Is the V-Axis the Only Direction?}
\label{sec:counterfactual}

To quantify precisely what the V-axis captures, we follow the
activation-ablation protocol of
\cite{arditi2024refusal, park2024linear}. At each of L24, L27, L28 we
(i)~compute the 5-fold cross-validated AUC of a full-rank $L_2$
logistic regression on the hidden features, (ii)~project the V-axis
out of the features (orthogonal complement), and (iii)~refit the same
full-rank LogReg on the ablated features.

\begin{figure}[t]
\centering
\includegraphics[width=0.85\textwidth]{fig_counterfactual}
\caption{V-axis ablation at L24/L27/L28 on SST-2. The blue bar is the
median-thresh 1-D V-axis probe; green is full-rank supervised LogReg
on the hidden state; orange is the same LogReg after projecting the
V-axis out. The change from green to orange is at most $0.003$ AUC,
indicating that valence is distributed across many directions and the
V-axis we extract is a \emph{summary}, not a bottleneck.}
\label{fig:counterfactual}
\end{figure}

\begin{table}[h]
\centering
\caption{Counterfactual V-axis ablation on SST-2 (5-fold CV AUC,
Qwen2.5-1.5B). The full-rank LogReg baseline is the ceiling reachable by
a supervised linear probe on the full hidden state; ``ablated'' is the
same LogReg after the 1-D V-axis is projected out.}
\label{tab:counterfactual}
\begin{tabular}{lccc}
\toprule
Layer & 1-D V-axis (ours) & Full-rank LogReg & After V-axis ablation \\
\midrule
L24 & 0.675 & 0.916 & 0.914 \\
L27 & 0.857 & 0.911 & 0.908 \\
L28 & \textbf{0.868} & 0.911 & 0.910 \\
\bottomrule
\end{tabular}
\end{table}

The result is informative in both directions. A full-rank supervised
probe reaches AUC $0.911$; projecting our 1-D V-axis \emph{out} of the
features changes that by only $0.001$ ($0.910$). Valence is therefore
richly encoded across many directions --- it is not localised on a
single bottleneck axis --- and the right framing of \ours{} is a
compressed summary, not a causal handle. The same counterfactual
\emph{positively} sharpens the main result: a supervised full-rank
probe with access to labels and all $1{,}536$ dimensions beats us by
only $+0.04$ AUC, while we use zero labels and one direction.

\subsection{Unsupervised PCA Beats Supervised LDA Out-of-Distribution}
\label{sec:pca_vs_lda}

Two ways to extract a concept direction from the story embeddings:
(i)~PCA of the 9 centroids (our method, unsupervised);
(ii)~supervised linear discriminant analysis (LDA) on the per-story
vectors (450 samples, 9 classes).

\begin{table}[h]
\centering
\caption{PCA beats LDA for OOD valence transfer.}
\label{tab:pca_vs_lda}
\begin{tabular}{lcc}
\toprule
Method & EmoBank V $|r|$ \\
\midrule
PCA (unsupervised, 9 centroids) & \textbf{0.498} \\
LDA (supervised, 450 per-story vectors) & 0.119 \\
\bottomrule
\end{tabular}
\end{table}

LDA overfits the 450 training stories and produces a direction that
does not transfer out-of-distribution; unsupervised PCA of the nine
centroids is $4\times$ more correlated with EmoBank. For
concept-direction extraction, unsupervised prototype geometry beats
supervised discriminant geometry whenever downstream data drifts from
the extraction set.

\subsection{Speed Benchmark}

\begin{table}[h]
\centering
\caption{Probing vs.\ prompting inference speed (per SST-2 sample,
Qwen2.5-1.5B on V100).}
\label{tab:speed}
\begin{tabular}{lc}
\toprule
Method & ms/sample \\
\midrule
\ours (batched forward pass + projection) & 16.07 \\
Direct prompting (serial, greedy decode) & 123.58 \\
\midrule
\rowcolor{blond}
\textbf{Speedup} & \textbf{\speedup{}} \\
\bottomrule
\end{tabular}
\end{table}

Probing is \speedup{} faster than prompting. The advantage grows with batch
size (probing batches naturally; autoregressive prompting is serial).

\subsection{Baseline Comparison on Standard Sentiment Tasks}

\begin{table}[h]
\centering
\caption{Comprehensive baseline comparison. Our method sits between direct
prompting and SBERT/random baselines.}
\label{tab:allbaselines}
\begin{tabular}{lccc}
\toprule
Method & SST-2 AUC & IMDB AUC & Hu \& Liu $|r|$ \\
\midrule
Random projection (mean of 10) & 0.509 & 0.508 & 0.12 \\
SBERT (MiniLM-L6, same pipeline) & 0.713 & 0.722 & 0.52 \\
\rowcolor{blond}
\textbf{\ours (Qwen2.5-1.5B, zero labels)} & \textbf{0.868} & \textbf{0.895} & \textbf{0.76} \\
Supervised LR ($N=5000$) & 0.914 & -- & -- \\
Direct prompting (greedy) & 0.928 & 0.860 & 0.840 \\
\bottomrule
\end{tabular}
\end{table}

\subsection{Scope of the Method}
\label{sec:scope}

The preceding experiments also delineate where \ours{} stops being the
right tool. We quantify four complementary bounds.

\paragraph{Concept class: global-polarity concepts.}
Repeating the 9-centroid recipe with 15 toxic and 15 safe stories produces
a T-axis that reaches AUC $0.59$ on Jigsaw Civil Comments (near chance).
Toxicity depends on specific lexical triggers whose presence is not
summarised by a single global direction learned from non-toxic
narrative stories. The method works when the concept has a clean
global-polarity linear manifestation; see Supplementary for the
toxicity trace.

\paragraph{Arousal bounds the regime.}
Applying the same PCA recipe to the second principal component does not
recover a usable arousal axis: across every layer and every model we
tested, EmoBank arousal correlation stays at $|r|\leq 0.16$. Combined
with Section~\ref{sec:va_asymmetry} this bounds the operating regime of
our method to \emph{valence-like global concepts aligned with the linear
representation hypothesis} \cite{park2024linear}. Arousal either lives on
a non-linear geometry in the residual stream or is encoded non-locally.
We treat this as a falsifiable scope statement, not a weakness: the
method's regime is the set of concepts whose linear directions survive
the late-layer re-linearisation.

\paragraph{Supervised and prompt baselines dominate on raw accuracy.}
Direct few-shot prompting reaches $92.78\%$ on SST-2, $15$ points above
our median-threshold probe, and a full-rank supervised logistic
regression on the same hidden states reaches AUC $0.911$ ($+0.043$).
\ours{} sits below these baselines on raw accuracy \emph{by
construction}: it is a compressed, zero-label, 1-D probe, not a
full-capacity classifier. The value is in data-efficiency, multilingual
transfer, interpretability, and continuous scores.

\paragraph{Cross-LLM ensembling does not help.}
A sign-aligned mean projection across Qwen-0.5B, Qwen-1.5B, Pythia-1.4B,
and Bloom-560m underperforms the best single LLM (Qwen-1.5B) by
$0.02$--$0.04$ AUC on SST-2 / IMDB / Yelp. Weaker extractors pull the
ensemble toward a less-aligned direction. The practical recommendation
is to pick the best single extractor and use its axis directly.
